\section{Optimizer Families}
\label{sec:annex-optimizers}

\subsection{Memory and Complexity Comparison}

Table~\ref{tab:optimizer-complexity} summarizes the memory overhead (number of state buffers per parameter), per-step computational complexity, and key hyperparameters for all supported optimizers. Memory overhead is expressed in terms of the number of state tensors maintained per model parameter, where each tensor has the same shape as the parameter.

\begin{longtable}{llcl}
\toprule
\textbf{Optimizer} & \textbf{State Buffers} & \textbf{Per-Step Cost} & \textbf{Key Hyperparams} \\
\midrule
\endfirsthead
\toprule
\textbf{Optimizer} & \textbf{State Buffers} & \textbf{Per-Step Cost} & \textbf{Key Hyperparams} \\
\midrule
\endhead
\midrule
\multicolumn{4}{r}{\textit{continued on next page}} \\
\endfoot
\bottomrule
\caption{Memory and complexity comparison of all supported optimizers. $n$ = number of parameters, $d$ = tensor dimensionality, $r$ = low-rank dimension, $k$ = Newton--Schulz iteration count. State buffers lists the number of optimizer state tensors maintained per parameter group. Per-step cost focuses on the dominant additional cost beyond the gradient computation itself.}
\label{tab:optimizer-complexity}
\endlastfoot
SGD+Momentum & 1 ($m$) & $\mathcal{O}(n)$ & lr, momentum, wd \\
AdamW & 2 ($m$, $v$) & $\mathcal{O}(n)$ & lr, $\beta_1$, $\beta_2$, eps, wd \\
RAdam & 2 ($m$, $v$) & $\mathcal{O}(n)$ & lr, $\beta_1$, $\beta_2$, eps, wd \\
Adan & 3 ($m$, $v$, $s$) & $\mathcal{O}(n)$ & lr, $\beta_1$, $\beta_2$, $\beta_3$, eps, wd \\
ADOPT & 2 ($m$, $v$) & $\mathcal{O}(n)$ & lr, $\beta_1$, $\beta_2$, eps, wd \\
AdEMAMix & 3 ($m_1$, $m_2$, $v$) & $\mathcal{O}(n)$ & lr, $\beta_1$, $\beta_2$, $\beta_3$, eps, wd \\
MARS & 3 ($m$, $v$, $z$) & $\mathcal{O}(n)$ & lr, $\beta_1$, $\beta_2$, eps, wd, $\gamma$ \\
Cautious AdamW & 2 ($m$, $v$) & $\mathcal{O}(n)$ & lr, $\beta_1$, $\beta_2$, eps, wd \\
LAMB & 2 ($m$, $v$) & $\mathcal{O}(n)$ & lr, $\beta_1$, $\beta_2$, eps, wd \\
Schedule-Free & 3 ($z$, $x$, $n$) & $\mathcal{O}(n)$ & lr, $\beta_1$, $\beta_2$, wd \\
Adafactor & 1--2 (row/col) & $\mathcal{O}(n)$ & lr, $\beta_1$, $\beta_2$, eps, wd, factored \\
GaLore & 2 ($m$, $v$) + SVD/proj & $\mathcal{O}(nr)$ & lr, rank $r$, $\beta_1$, $\beta_2$, eps, wd \\
Prodigy & 3 ($m$, $v$, $d$) & $\mathcal{O}(n)$ & $\beta_1$, $\beta_2$, eps, wd, $d_0$ \\
Lion & 1 ($m$) & $\mathcal{O}(n)$ & lr, $\beta_1$, $\beta_2$, wd \\
Shampoo & $2d$ ($L_i$, $R_i$) & $\mathcal{O}(n^{1+2/d})$ & lr, eps, wd, matrix eps \\
SOAP & $2d$ + 2 ($m$, $v$) & $\mathcal{O}(n^{1+1/d})$ & lr, $\beta_1$, $\beta_2$, eps, wd, shampoo eps \\
Sophia & 3 ($m$, $v$, $h$) & $\mathcal{O}(n)$ & lr, $\beta_1$, $\beta_2$, eps, wd, $k$ \\
Muon & 2 ($m$, $v$) & $\mathcal{O}(n \cdot k)$ & lr, momentum, wd, NS steps $k$ \\
Turbo-Muon & 2 ($m$, $v$) & $\mathcal{O}(n \cdot k)$ & lr, momentum, wd, NS steps $k$ \\
APOLLO & 2 low-rank ($m_R$, $v_R$) + proj & $\mathcal{O}(nr)$ & lr, rank $r$, update gap, scale, betas, eps, wd \\
APOLLO-Mini & 2 rank-1 ($m_R$, $v_R$) + proj & $\mathcal{O}(n)$ & lr, update gap, scale, betas, eps, wd \\
Q-APOLLO & 2 quantized low-rank + proj & $\mathcal{O}(nr)$ & lr, rank $r$, update gap, scale, quant bits, betas, eps, wd \\
Anon & 3 ($m$, $v$, fixed\_lr) & $\mathcal{O}(n)$ & lr, $\beta_1$, $\beta_2$, eps, wd, $\gamma$ \\
\end{longtable}

\subsection{The Evolution of Optimization in Neural Networks}
The optimizer survey frames transformer optimization as a response to three structural pressures: non-convex loss landscapes, severe curvature heterogeneity across parameter blocks, and the memory cost of storing optimizer state for very large models. The report argues that the field has diverged into several trajectories: adaptive first-order baselines, variance-reduction methods, memory-efficient methods, structured second-order preconditioners, schedule-free methods, and orthogonality-oriented updates.

\subsection{Standard Baseline and Adaptive Optimizers}
\paragraph{SGD with Momentum.}
The classical update accumulates a momentum buffer and then applies a fixed learning rate. At each iteration $t$, the algorithm maintains an exponential moving average $m_t$ of past gradients:
\begin{align}
m_t &= \beta m_{t-1} + g_t \\
\theta_{t+1} &= \theta_t - \eta m_t
\end{align}
where $g_t = \nabla f(\theta_t)$, $\beta \in [0.8, 0.99]$ controls the momentum decay, and $\eta$ is the fixed learning rate. The momentum term acts as velocity: it accelerates movement in consistent directions while dampening oscillations in variable directions. Its strengths are low memory overhead (single momentum buffer per parameter) and strong generalization when tuned carefully. Its main weakness in transformer workloads is poor robustness to heterogeneous curvature (different Hessian spectra across parameter groups) and strong dependence on learning-rate schedules.

\begin{algorithm}[H]
\caption{SGD with Momentum}
\label{alg:sgd-momentum}
\begin{algorithmic}[1]
\Require Initial parameters $\theta_0$, learning rate $\eta$, momentum coefficient $\beta$, weight decay $\lambda$
\Ensure Updated parameters $\theta_T$
\State Initialize: momentum buffer $m \gets 0$
\For{$t = 0, 1, 2, \ldots, T-1$}
    \State Compute gradient: $g_t \gets \nabla f(\theta_t)$
    \If{weight\_decay $> 0$}
        \State $g_t \gets g_t + \lambda \theta_t$ \Comment{L2 regularization}
    \EndIf
    \State Update momentum: $m \gets \beta \cdot m + g_t$
    \State Update parameters: $\theta_{t+1} \gets \theta_t - \eta \cdot m$
\EndFor
\Return $\theta_T$
\end{algorithmic}
\end{algorithm}

\paragraph{Adam and AdamW.}
Adam tracks exponential moving averages of both first moment (mean) and second moment (uncentered variance) to achieve element-wise adaptive learning rates. The update rule is:
\begin{align}
m_t &= \beta_1 m_{t-1} + (1-\beta_1) g_t \\
v_t &= \beta_2 v_{t-1} + (1-\beta_2) g_t^2 \\
\hat{m}_t &= \frac{m_t}{1-\beta_1^t}, \quad \hat{v}_t = \frac{v_t}{1-\beta_2^t} \quad \text{(bias correction)} \\
\theta_{t+1} &= \theta_t - \eta \left(\frac{\hat{m}_t}{\sqrt{\hat{v}_t}+\epsilon}\right)
\end{align}
AdamW introduces a crucial modification: weight decay is applied directly to parameters (decoupled from gradients): $\theta_t \gets \theta_t(1-\eta\lambda)$ rather than adding $\lambda\theta_t$ to the gradient. This decoupling prevents the adaptive scaling from interfering with regularization strength. The report treats AdamW as the practical baseline for transformer fine-tuning because it converges quickly and is relatively forgiving to hyperparameter variation. The tradeoff is memory cost: both moment tensors must be stored for every parameter, doubling optimizer-state memory relative to SGD.

\begin{algorithm}[H]
\caption{AdamW (Adam with Decoupled Weight Decay)}
\label{alg:adamw}
\begin{algorithmic}[1]
\Require Initial parameters $\theta_0$, learning rate $\eta$, exponential decay rates $\beta_1, \beta_2 \in [0,1)$
\Require Momentum constant $\epsilon$, weight decay $\lambda$
\Ensure Updated parameters $\theta_T$
\State Initialize: first moment $m \gets 0$, second moment $v \gets 0$, step counter $t \gets 0$
\For{$t = 1, 2, \ldots, T$}
    \State Compute gradient: $g_t \gets \nabla f(\theta_{t-1})$
    \State Weight decay (decoupled): $\theta_{t-1} \gets \theta_{t-1}(1 - \eta\lambda)$
    \State Update moments: $m \gets \beta_1 m + (1-\beta_1) g_t$
    \State $v \gets \beta_2 v + (1-\beta_2) g_t^2$
    \State Bias correction: $\hat{m} \gets m / (1 - \beta_1^t)$, $\hat{v} \gets v / (1 - \beta_2^t)$
    \State Update parameters: $\theta_t \gets \theta_{t-1} - \eta (\hat{m} / (\sqrt{\hat{v}} + \epsilon))$
\EndFor
\Return $\theta_T$
\end{algorithmic}
\end{algorithm}

\paragraph{RAdam (Rectified Adam).}
RAdam addresses Adam's instability in early training by dynamically rectifying the adaptive learning rate. The key observation is that the second moment $v_t$ has very high variance in the first few steps, causing unreliable adaptive scaling. RAdam computes the effective simple moving average (SMA) window length:
\[\rho_t = \rho_{\infty} - \frac{2t\beta_2^t}{1-\beta_2^t}, \quad \text{where} \quad \rho_{\infty} = \frac{2}{1-\beta_2}-1\]
When $\rho_t > 4$ (sufficient samples for variance estimation), RAdam applies adaptive scaling with a rectification term:
\[r_t = \sqrt{\frac{(\rho_t-4)(\rho_t-2)\rho_{\infty}}{(\rho_{\infty}-4)(\rho_{\infty}-2)\rho_t}}, \quad \theta_{t+1} = \theta_t - \eta r_t \frac{\hat{m}_t}{\sqrt{\hat{v}_t}+\epsilon}\]
When $\rho_t \leq 4$, RAdam falls back to SGD with momentum: $\theta_{t+1} = \theta_t - \eta\hat{m}_t$. This graceful transition eliminates the need for manual learning-rate warmup schedules and improves robustness.

\begin{algorithm}[H]
\caption{RAdam (Rectified Adam)}
\label{alg:radam}
\begin{algorithmic}[1]
\Require Initial parameters $\theta_0$, learning rate $\eta$, $\beta_1, \beta_2$, weight decay $\lambda$
\Ensure Updated parameters $\theta_T$
\State Initialize: $m \gets 0$, $v \gets 0$, $t \gets 0$
\State Compute: $\rho_\infty \gets 2/(1-\beta_2) - 1$
\For{$t = 1, 2, \ldots, T$}
    \State $g_t \gets \nabla f(\theta_{t-1})$
    \If{weight\_decay $> 0$}
        \State $\theta_{t-1} \gets \theta_{t-1}(1 - \eta\lambda)$
    \EndIf
    \State $m \gets \beta_1 m + (1-\beta_1) g_t$
    \State $v \gets \beta_2 v + (1-\beta_2) g_t^2$
    \State $\hat{m} \gets m/(1-\beta_1^t)$, $\hat{v} \gets v/(1-\beta_2^t)$
    \State $\rho_t \gets \rho_\infty - 2t \beta_2^t / (1 - \beta_2^t)$
    \If{$\rho_t > 4$}
        \State $r_t \gets \sqrt{((\rho_t-4)(\rho_t-2)\rho_\infty) / ((\rho_\infty-4)(\rho_\infty-2)\rho_t)}$
        \State $\theta_t \gets \theta_{t-1} - \eta r_t \hat{m} / (\sqrt{\hat{v}} + \epsilon)$
    \Else
        \State $\theta_t \gets \theta_{t-1} - \eta \hat{m}$ \Comment{SGD with momentum}
    \EndIf
\EndFor
\Return $\theta_T$
\end{algorithmic}
\end{algorithm}

\subsection{Advanced Momentum and Variance Reduction (2024--2025)}
\paragraph{Adan (Adaptive Nesterov Momentum).}
Adan reformulates Nesterov acceleration without the extra gradient computation required by classical Nesterov SGD. The algorithm maintains three momentum buffers:
\begin{align}
m_t &= (1-\beta_1)m_{t-1} + \beta_1 g_t \quad\text{(first moment)} \\
v_t &= (1-\beta_2)v_{t-1} + \beta_2(g_t - g_{t-1}) \quad\text{(velocity/gradient difference)} \\
n_t &= (1-\beta_3)n_{t-1} + \beta_3[g_t + (1-\beta_1)(g_t-g_{t-1})]^2 \quad\text{(Nesterov second moment)}
\end{align}
The \textbf{Nesterov Momentum Estimation} (NME) term $\bar{g}_t = g_t + (1-\beta_1)(g_t-g_{t-1})$ estimates the gradient at a future position without evaluating it. The update combines momentum with velocity:
\[\bar{m}_t = m_t + (1-\beta_1)v_t, \quad \theta_{t+1} = \theta_t - \eta\frac{\bar{m}_t}{\sqrt{n_t}+\epsilon}\]
Adan achieves fast convergence across diverse architectures (CNNs, GANs, Transformers) through this acceleration. The cost is maintaining three momentum-like buffers, increasing memory overhead relative to Adam.

\begin{algorithm}[H]
\caption{Adan (Adaptive Nesterov Momentum)}
\label{alg:adan}
\begin{algorithmic}[1]
\Require Initial parameters $\theta_0$, learning rate $\eta$, $\beta_1, \beta_2, \beta_3 \in (0,1)$
\Ensure Updated parameters $\theta_T$
\State Initialize: $m \gets 0$, $v \gets 0$, $n \gets 0$, $g_{-1} \gets 0$ (previous gradient)
\For{$t = 1, 2, \ldots, T$}
    \State $g_t \gets \nabla f(\theta_{t-1})$
    \State $m \gets (1-\beta_1) m + \beta_1 g_t$
    \State $\Delta g_t \gets g_t - g_{t-1}$
    \State $v \gets (1-\beta_2) v + \beta_2 \Delta g_t$
    \State Nesterov estimation: $\bar{g}_t \gets g_t + (1-\beta_1) \Delta g_t$
    \State $n \gets (1-\beta_3) n + \beta_3 \bar{g}_t^2$
    \State Combined momentum: $\bar{m} \gets m + (1-\beta_1) v$
    \State $\theta_t \gets \theta_{t-1} - \eta (\bar{m} / (\sqrt{n} + \epsilon))$
    \State $g_{t-1} \gets g_t$
\EndFor
\Return $\theta_T$
\end{algorithmic}
\end{algorithm}

\paragraph{ADOPT (Adam with Optimal Pruned Tuning).}
ADOPT fixes a fundamental theoretical issue in Adam: the gradient appears in both the first moment and second moment estimates, creating circularity. ADOPT \textbf{decouples} by using the previous step's second moment for the denominator:
\begin{align}
v_t &= \beta_2 v_{t-1} + (1-\beta_2)g_t^2 \\
m_t &= \beta_1 m_{t-1} + (1-\beta_1)\frac{g_t}{\sqrt{v_{t-1}}+\epsilon} \quad\text{(uses $v_{t-1}$, not $v_t$)} \\
\theta_{t+1} &= \theta_t - \eta m_t
\end{align}
This simple reordering achieves the optimal convergence rate $O(1/\sqrt{T})$ with \textbf{any} choice of $\beta_2 \in (0,1)$, without bounded-noise assumptions. The practical consequence is that ADOPT is a drop-in replacement for Adam with stronger theoretical guarantees and comparable or superior empirical performance across vision, NLP, RL, and generative modeling domains.

\begin{algorithm}[H]
\caption{ADOPT (Adam with Optimal Pruned Tuning)}
\label{alg:adopt}
\begin{algorithmic}[1]
\Require Initial parameters $\theta_0$, learning rate $\eta$, $\beta_1, \beta_2$, tolerance $\epsilon$
\Ensure Updated parameters $\theta_T$
\State Initialize: $m \gets 0$, $v \gets 0$
\For{$t = 1, 2, \ldots, T$}
    \State $g_t \gets \nabla f(\theta_{t-1})$
    \If{weight\_decay $> 0$}
        \State $\theta_{t-1} \gets \theta_{t-1}(1 - \eta\lambda)$
    \EndIf
    \State Use \textbf{previous} second moment: $\text{denom} \gets \sqrt{\max(v, \epsilon)} + \epsilon$
    \State Update first moment using previous variance: $m \gets \beta_1 m + (1-\beta_1)(g_t / \text{denom})$
    \State Update second moment with \textbf{current} gradient: $v \gets \beta_2 v + (1-\beta_2) g_t^2$
    \State $\theta_t \gets \theta_{t-1} - \eta m$
\EndFor
\Return $\theta_T$
\end{algorithmic}
\end{algorithm}

\paragraph{AdEMAMix (Exponential Moving Average Mixture).}
AdEMAMix replaces Adam's single EMA of gradients with a \textbf{mixture of two EMAs}: one fast-decaying and one slow-decaying. This addresses the observation that gradients remain informative over tens of thousands of steps, not just hundreds:
\begin{align}
m_{1,t} &= \beta_1 m_{1,t-1} + (1-\beta_1)g_t \quad\text{(fast EMA, $\beta_1 \approx 0.9$)} \\
m_{2,t} &= \beta_3 m_{2,t-1} + (1-\beta_3)g_t \quad\text{(slow EMA, $\beta_3 \approx 0.9999$)} \\
v_t &= \beta_2 v_{t-1} + (1-\beta_2)g_t^2 \\
m_t &= m_{1,t} + \alpha_t m_{2,t} \quad\text{(mixture with scheduled weight $\alpha_t$)} \\
\theta_{t+1} &= \theta_t - \eta\frac{m_t}{\sqrt{v_t}+\epsilon}
\end{align}
The mixture weight $\alpha_t$ typically increases during training, allowing the fast EMA to provide immediate adaptation while the slow EMA accumulates long-range gradient correlations. Empirically, a 1.3B model on 101B tokens achieves similar loss to AdamW on 197B tokens (95\% data efficiency gain), suggesting the slow EMA significantly reduces model forgetting.

\begin{algorithm}[H]
\caption{AdEMAMix (Exponential Moving Average Mixture)}
\label{alg:ademamix}
\begin{algorithmic}[1]
\Require Initial parameters $\theta_0$, learning rate $\eta$, $\beta_1$ (fast), $\beta_3$ (slow), $\beta_2$ (second moment)
\Require Mixture weight $\alpha_t$ (typically increasing), tolerance $\epsilon$
\Ensure Updated parameters $\theta_T$
\State Initialize: $m_1 \gets 0$ (fast EMA), $m_2 \gets 0$ (slow EMA), $v \gets 0$
\For{$t = 1, 2, \ldots, T$}
    \State $g_t \gets \nabla f(\theta_{t-1})$
    \If{weight\_decay $> 0$}
        \State $\theta_{t-1} \gets \theta_{t-1}(1 - \eta\lambda)$
    \EndIf
    \State $m_1 \gets \beta_1 m_1 + (1-\beta_1) g_t$ \Comment{Fast EMA}
    \State $m_2 \gets \beta_3 m_2 + (1-\beta_3) g_t$ \Comment{Slow EMA}
    \State $v \gets \beta_2 v + (1-\beta_2) g_t^2$
    \State $m_t \gets m_1 + \alpha_t m_2$ \Comment{Mixture}
    \State $\theta_t \gets \theta_{t-1} - \eta (m_t / (\sqrt{v} + \epsilon))$
\EndFor
\Return $\theta_T$
\end{algorithmic}
\end{algorithm}

\paragraph{MARS (Make Adaptive learning Rates Shine).}
MARS combines preconditioned gradient methods (e.g., AdamW) with \textbf{variance reduction} via scaled stochastic recursive momentum. The core innovation is a variance-reduced gradient estimate:
\begin{align}
c_t &= g_t + \gamma(c_{t-1} - g_{t-1}) \quad\text{(SVRG-style recursive estimate)} \\
m_t &= \beta_1 m_{t-1} + (1-\beta_1)c_t \\
v_t &= \beta_2 v_{t-1} + (1-\beta_2)c_t^2 \\
\theta_{t+1} &= \theta_t - \eta\frac{m_t}{\sqrt{v_t}+\epsilon}
\end{align}
where $\gamma \in [0.01, 0.1]$ controls how much historical gradient information is retained. The variance-reduced gradient $c_t$ acts as an implicit noise filter, amplifying consistent signal directions and dampening contradictory noise. MARS-AdamW consistently outperforms bare AdamW by significant margins on GPT-2 pretraining, suggesting that variance reduction substantially improves convergence in mini-batch stochastic training.

\begin{algorithm}[H]
\caption{MARS (Make Adaptive learning Rates Shine)}
\label{alg:mars}
\begin{algorithmic}[1]
\Require Initial parameters $\theta_0$, learning rate $\eta$, $\beta_1, \beta_2$, weight decay $\lambda$
\Require Variance reduction coefficient $\gamma \in [0.01, 0.1]$
\Ensure Updated parameters $\theta_T$
\State Initialize: $m \gets 0$, $v \gets 0$, $c \gets 0$ (variance-reduced gradient), $g_{-1} \gets 0$
\For{$t = 1, 2, \ldots, T$}
    \State $g_t \gets \nabla f(\theta_{t-1})$
    \State Variance reduction: $c \gets g_t + \gamma(c - g_{t-1})$
    \State $m \gets \beta_1 m + (1-\beta_1) c$
    \State $v \gets \beta_2 v + (1-\beta_2) c^2$
    \If{weight\_decay $> 0$}
        \State $\theta_{t-1} \gets \theta_{t-1}(1 - \eta\lambda)$
    \EndIf
    \State $\theta_t \gets \theta_{t-1} - \eta (m / (\sqrt{v} + \epsilon))$
    \State $g_{t-1} \gets g_t$
\EndFor
\Return $\theta_T$
\end{algorithmic}
\end{algorithm}

\paragraph{Cautious Optimizers (Cautious AdamW, Cautious Lion).}
The Cautious framework applies a \textbf{one-line modification} to any momentum-based optimizer: mask the update so that only dimensions where momentum and gradient directions agree are applied:
\begin{align}
\text{mask}_i &= \begin{cases} 1 & \text{if } m_t[i] \cdot g_t[i] > 0 \quad\text{(agreement)} \\ 0 & \text{otherwise} \end{cases} \\
u_t &= \left(\frac{m_t}{\sqrt{v_t}+\epsilon}\right) \odot \text{mask} \quad\text{(element-wise masking)} \\
\theta_{t+1} &= \theta_t - \eta u_t
\end{align}
The intuition: momentum $m_t$ estimates the gradient direction from history; the current gradient $g_t$ is instantaneous signal. When both agree, the optimizer is confident and should update aggressively. When they disagree, the optimizer is conflicted---the historical trend points toward a region that may have been good before, but current evidence contradicts it. By masking conflicted dimensions, Cautious becomes more conservative and avoids corrupted steps. Empirically, this simple masking achieves up to \textbf{1.47$\times$ speedup} on Llama and MAE pretraining while preserving convergence guarantees.

\begin{algorithm}[H]
\caption{Cautious AdamW (Consensus-based Update Masking)}
\label{alg:cautious}
\begin{algorithmic}[1]
\Require Initial parameters $\theta_0$, learning rate $\eta$, $\beta_1, \beta_2$, weight decay $\lambda$
\Ensure Updated parameters $\theta_T$
\State Initialize: $m \gets 0$, $v \gets 0$
\For{$t = 1, 2, \ldots, T$}
    \State $g_t \gets \nabla f(\theta_{t-1})$
    \State $m \gets \beta_1 m + (1-\beta_1) g_t$
    \State $v \gets \beta_2 v + (1-\beta_2) g_t^2$
    \State Consensus mask: $\text{mask}_i \gets 1$ if $m[i] \cdot g_t[i] > 0$, else $0$ \Comment{Agreement}
    \State Base Adam update: $u \gets m / (\sqrt{v} + \epsilon)$
    \State Apply mask: $u_{\text{masked}} \gets u \odot \text{mask}$ \Comment{Element-wise}
    \If{weight\_decay $> 0$}
        \State $\theta_{t-1} \gets \theta_{t-1}(1 - \eta\lambda)$
    \EndIf
    \State $\theta_t \gets \theta_{t-1} - \eta u_{\text{masked}}$
\EndFor
\Return $\theta_T$
\end{algorithmic}
\end{algorithm}

\subsection{Large-Batch, Memory-Efficient, and Parameter-Free Optimizers}
\paragraph{LAMB (Layer-wise Adaptive Moments optimizer for Batch training).}
LAMB extends Adam with \textbf{layer-wise adaptive rate scaling} (inspired by LARS), enabling stable training with extreme batch sizes (e.g., 64K on BERT):
\begin{align}
m_t^L &= \beta_1 m_{t-1}^L + (1-\beta_1)g_t^L \quad\text{(layer-wise first moment)} \\
v_t^L &= \beta_2 v_{t-1}^L + (1-\beta_2)(g_t^L)^2 \\
u_{\text{adam}}^L &= \frac{m_t^L}{\sqrt{v_t^L}+\epsilon} + \lambda\theta_{t-1}^L \quad\text{(base adaptive step with decay)} \\
\phi^L &= \frac{\|\theta_{t-1}^L\|_2}{\|u_{\text{adam}}^L\|_2} \quad\text{(trust ratio: layer normalization)} \\
\theta_t^L &= \theta_{t-1}^L - \eta \cdot \phi^L \cdot u_{\text{adam}}^L
\end{align}
The \textbf{trust ratio} $\phi^L = \|\theta^L\|_2 / \|u_{\text{adam}}^L\|_2$ normalizes the effective update step relative to weight magnitude. In very large batches, stochastic gradient noise can exceed signal, destabilizing layer-wise learning rates. By scaling updates proportionally to weight norms, LAMB ensures relative changes rather than absolute ones, preventing small confident updates in large vectors from dominating. LAMB preserves small-batch generalization benefits while enabling efficient large-batch training.

\begin{algorithm}[H]
\caption{LAMB (Layer-wise Adaptive Moments optimizer for Batch training)}
\label{alg:lamb}
\begin{algorithmic}[1]
\Require Initial parameters $\theta_0$, learning rate $\eta$, $\beta_1, \beta_2$, weight decay $\lambda$
\Ensure Updated parameters $\theta_T$
\State Initialize: For each layer $L$: $m^L \gets 0$, $v^L \gets 0$
\For{$t = 1, 2, \ldots, T$}
    \For{each layer $L$}
        \State $g_t^L \gets \nabla f(\theta_t^L)$
        \State $m^L \gets \beta_1 m^L + (1-\beta_1) g_t^L$
        \State $v^L \gets \beta_2 v^L + (1-\beta_2) (g_t^L)^2$
        \State Base Adam: $u_{\text{adam}}^L \gets m^L / (\sqrt{v^L} + \epsilon)$
        \If{weight\_decay $> 0$}
            \State $u_{\text{adam}}^L \gets u_{\text{adam}}^L + \lambda \theta_{t-1}^L$
        \EndIf
        \State Trust ratio: $\phi^L \gets \|\theta_{t-1}^L\|_2 / \|u_{\text{adam}}^L\|_2$ if both nonzero, else $1$
        \State $\theta_t^L \gets \theta_{t-1}^L - \eta \phi^L u_{\text{adam}}^L$
    \EndFor
\EndFor
\Return $\theta_T$
\end{algorithmic}
\end{algorithm}

\paragraph{Schedule-Free AdamW.}
Schedule-free methods remove explicit scheduler design from the optimization recipe. The algorithm maintains two parameter streams: $z_t$ (exploration) and $x_t$ (smooth average):
\begin{align}
v_t &= \beta_2 v_{t-1} + (1-\beta_2)g_t^2 \quad\text{(variance)} \\
z_t &= z_{t-1}(1-\eta\lambda) - \eta\frac{g_t}{\sqrt{v_t}+\epsilon} \quad\text{(adaptive step with decay)} \\
c_t &= \frac{1}{t+1} \quad\text{(averaging coefficient, decreases over time)} \\
x_t &= (1-c_t)x_{t-1} + c_t z_t \quad\text{(iterate averaging)} \\
y_t &= (1-\beta)z_t + \beta x_t \quad\text{(interpolation for evaluation point)}
\end{align}
The averaging coefficient $c_t = 1/(t+1)$ implements an implicit learning-rate schedule without explicitly specifying total steps $T$. This framework unifies scheduling and iterate averaging: the algorithm explores via $z_t$ while accumulating stable direction via $x_t$. The evaluation point $y_t$ (used for gradient computation) interpolates between exploration and stability. Schedule-Free achieves state-of-the-art convergence across convex optimization, large-scale deep learning, and reinforcement learning, while removing a major hyperparameter.

\begin{algorithm}[H]
\caption{Schedule-Free AdamW}
\label{alg:schedulefree}
\begin{algorithmic}[1]
\Require Initial parameters $\theta_0$, learning rate $\eta$ (fixed), $\beta_1, \beta_2$, weight decay $\lambda$
\Ensure Updated parameters $\theta_T$ or averaging $x_T$
\State Initialize: $z \gets \theta_0$ (exploration), $x \gets \theta_0$ (average), $v \gets 0$, $t \gets 0$
\For{$t = 1, 2, \ldots, T$}
    \State $g_t \gets \nabla f(y_{t-1})$ (gradient at interpolation point)
    \State $v \gets \beta_2 v + (1-\beta_2) g_t^2$ \Comment{Variance}
    \State Weight decay on $z$: $z \gets z(1 - \eta\lambda)$
    \State $z \gets z - \eta g_t / (\sqrt{v} + \epsilon)$ \Comment{Adaptive step on raw}
    \State $c_t \gets 1 / (t+1)$ \Comment{Averaging coefficient}
    \State $x \gets (1-c_t) x + c_t z$ \Comment{Iterate averaging}
    \State $y_t \gets (1-\beta_1) z + \beta_1 x$ \Comment{Interpolation for next eval}
\EndFor
\Return $x_T$ or $y_T$
\end{algorithmic}
\end{algorithm}

\paragraph{Adafactor.}
Adafactor reduces optimizer memory by \textbf{factorizing} second-moment statistics for matrix-shaped parameters, storing only row and column accumulators rather than dense variance states. For a gradient matrix $G_t \in \mathbb{R}^{m \times n}$:
\begin{align}
R_t &= \beta_2 R_{t-1} + (1-\beta_2)(G_t^2) \mathbf{1}_n^T \quad\text{(row variance)} \\
C_t &= \beta_2 C_{t-1} + (1-\beta_2) \mathbf{1}_m^\top (G_t^2) \quad\text{(column variance)} \\
\hat{V}_t &= \frac{R_t C_t}{\mathbf{1}_n^\top R_t} \quad\text{(reconstructed variance via outer product)} \\
U_t &= \frac{G_t}{\sqrt{\hat{V}_t}+\epsilon} \quad\text{(normalized adaptive step)} \\
\hat{U}_t &= \frac{U_t}{\max(1, \text{RMS}(U_t))} \quad\text{(stability clipping)}
\end{align}
Instead of storing $m \times n$ second-moment values, Adafactor stores only $m + n$ accumulators, reducing memory from $O(n_{\text{params}})$ to $O(\sqrt{n_{\text{params}}})$. This is most attractive when VRAM is dominated by optimizer state rather than activations. The tradeoff is reduced optimization expressiveness and potential instability on some tasks.

\begin{algorithm}[H]
\caption{Adafactor (Factorized Second-Moment Statistics)}
\label{alg:adafactor}
\begin{algorithmic}[1]
\Require Gradient matrix $G_t \in \mathbb{R}^{m \times n}$, learning rate $\eta$, $\beta_2 \approx 0.999$
\Ensure Updated parameters
\State Initialize: Row accumulators $R \gets \epsilon \mathbf{1}^T_m$, Column $C \gets \epsilon \mathbf{1}_n$
\For{$t = 1, 2, \ldots, T$}
    \State $G_t \gets \nabla f(\theta_t)$ (matrix gradient)
    \State $R \gets \beta_2 R + (1-\beta_2) (G_t^2) \mathbf{1}_n^T$ \Comment{Row inner products}
    \State $C \gets \beta_2 C + (1-\beta_2) \mathbf{1}_m^\top (G_t^2)$ \Comment{Column inner products}
    \State Reconstructed variance: $\hat{V} \gets (R \cdot C) / (\mathbf{1}_n^\top R)$ \Comment{Via outer product}
    \State Normalized adaptive step: $U_t \gets G_t / (\sqrt{\hat{V}} + \epsilon)$
    \State Per-row RMS clipping: $\hat{U}_t \gets U_t / \max(1, \text{RMS}(U_t))$
    \State $\theta_{t+1} \gets \theta_t - \eta \hat{U}_t$
\EndFor
\end{algorithmic}
\end{algorithm}

\paragraph{GaLore (Gradient Low-Rank Projection).}
GaLore projects 2D gradients into a low-rank subspace before optimization, reducing optimizer-state memory while maintaining adaptive step scaling. For a 2D parameter matrix $W \in \mathbb{R}^{m \times n}$:
\begin{align}
U, S, V &= \text{SVD}(G_t) \quad\text{(compute singular decomposition)} \\
P &\in \mathbb{R}^{n \times r} \quad\text{or} \quad P^\top \in \mathbb{R}^{r \times m} \quad\text{(select top-$r$ singular vectors)} \\
G_{\text{low}} &= P^\top G_t \quad\text{(project to low-rank space)} \\
\Delta_{\text{low}} &= \text{Adam}(G_{\text{low}}) \quad\text{(optimize in compressed space)} \\
\Delta &= P \Delta_{\text{low}} \quad\text{(reconstruct in original space)}
\end{align}
By projecting into a rank-$r$ subspace (typically $r \ll \min(m,n)$), optimizer state is reduced from $O(mn)$ to $O(r(m+n))$ for 2D parameters. This complementary memory-saving approach is especially relevant when the model is too large for full-rank optimizer state. GaLore works synergistically with other techniques and has shown strong empirical results on billion-parameter models.

\begin{algorithm}[H]
\caption{GaLore (Gradient Low-Rank Projection)}
\label{alg:galore}
\begin{algorithmic}[1]
\Require 2D parameter matrix $W \in \mathbb{R}^{m \times n}$, learning rate $\eta$, rank $r \ll \min(m,n)$
\Ensure Updated parameters
\State Initialize: Adam state in low-rank space (if rank-2 variant)
\For{$t = 1, 2, \ldots, T$}
    \State $G_t \gets \nabla f(W_t)$
    \If{$t \mod K = 1$}  \Comment{Periodic SVD}
        \State $U_t, S_t, V_t^\top \gets \text{SVD}(G_t)$ (full or thin)
        \State Select projection: $P \gets U_t[:, :r]$ or $P \gets V_t[:, :r]$
    \EndIf
    \State Project gradient: $G_{\text{low}} \gets P^\top G_t$ (or $G_t P^\top$ for right projection)
    \State Run Adam on low-rank: $\Delta_{\text{low}} \gets \text{Adam}(G_{\text{low}})$
    \State Reconstruct in original space: $\Delta \gets P \Delta_{\text{low}}$ (or $\Delta_{\text{low}} P^\top$)
    \State $W_t \gets W_t - \eta \Delta$
\EndFor
\end{algorithmic}
\end{algorithm}

\paragraph{APOLLO, APOLLO-Mini, and Q-APOLLO.}
The APOLLO family \citep{zhu2025apollosgdlikememoryadamwlevel} begins from the observation that AdamW's element-wise denominator can be \textbf{coarsened into a structured learning-rate update}. Rather than storing dense moments for every parameter entry, APOLLO projects a matrix gradient $G_t \in \mathbb{R}^{m \times n}$ into a compact random subspace and tracks Adam-style moments there:
\begin{align}
R_t &= P_t G_t \quad\text{or}\quad R_t = G_t P_t^\top \\
M_t^R &= \beta_1 M_{t-1}^R + (1-\beta_1)R_t \\
V_t^R &= \beta_2 V_{t-1}^R + (1-\beta_2)R_t^2 \\
\widetilde{R}_t &= \frac{M_t^R}{\sqrt{V_t^R}+\epsilon}
\end{align}
The projected state is \emph{not} expanded back into a dense low-rank update as in SVD-based methods. Instead, APOLLO estimates a structured scaling tensor $S_t$ in the original space. In the standard APOLLO variant, that scaling is \textbf{channel-wise}: each row or channel receives its own norm ratio. In APOLLO-Mini, the scaling is reduced to a single \textbf{tensor-wise} scalar, corresponding to the rank-1 extreme described in the paper. The resulting parameter update is therefore Adam-like in adaptation but much closer to SGD in state cost:
\[
W_t \gets (1-\eta\lambda)W_{t-1} - \eta\,\alpha\,(G_t \odot S_t)
\]
where $\alpha$ is the extra scale factor used to stabilize highly compressed variants.

The paper's practical contribution is twofold. First, APOLLO replaces expensive repeated SVD with periodically refreshed \textbf{Gaussian random projection}, so the systems burden is ordinary matrix multiplication rather than spectral decomposition. Second, the optimizer is unusually tolerant to extreme compression: even \textbf{APOLLO-Mini}, which keeps only rank-1 auxiliary state, remains competitive with or better than AdamW in the reported pre-training experiments while approaching SGD-level memory cost.

\begin{algorithm}[H]
\caption{APOLLO / APOLLO-Mini Structured Gradient Scaling}
\label{alg:apollo-family}
\begin{algorithmic}[1]
\Require Weight matrix $W \in \mathbb{R}^{m \times n}$ with $m \le n$, learning rate $\eta$, scale factor $\alpha$, decay rates $(\beta_1, \beta_2)$, weight decay $\lambda$, rank $r$, projection refresh interval $T$
\Ensure Updated parameters $W_T$
\State Initialize projected moments: $M^R \gets 0$, $V^R \gets 0$, step $t \gets 0$
\Repeat
    \State Compute gradient: $G_t \gets \nabla_W \phi(W_t)$
    \If{$t \bmod T = 0$}
        \State Sample Gaussian projector $P_t \sim \mathcal{N}(0, 1/r)$ with a fresh seed
    \EndIf
    \State Project gradient: $R_t \gets P_t G_t$ \Comment{or $G_t P_t^\top$ depending on layout}
    \State Update projected AdamW moments: $M_t^R, V_t^R \gets \text{AdamWState}(R_t; \beta_1, \beta_2)$
    \State Normalize projected state: $\widetilde{R}_t \gets M_t^R / (\sqrt{V_t^R}+\epsilon)$
    \If{APOLLO}
        \State $S_t \gets \operatorname{diag}(s_1^R, \ldots, s_m^R)$, where $s_i^R = \|\widetilde{R}_t[i,:]\|_2 / \|R_t[i,:]\|_2$
    \Else
        \State $S_t \gets s_t^R$, where $s_t^R = \|\widetilde{R}_t\|_2 / \|R_t\|_2$
    \EndIf
    \State Update weights: $W_t \gets (1-\eta\lambda)W_{t-1} - \eta\,\alpha\,(G_t \odot S_t)$
    \State $t \gets t+1$
\Until{convergence}
\end{algorithmic}
\end{algorithm}

\paragraph{APOLLO-Mini.}
APOLLO-Mini is the family member optimized for \textbf{extreme memory efficiency}. The paper motivates it by arguing that, in a rank-1 compact space, channel-wise scaling becomes too noisy, so the update is coarsened further to a single tensor-wise scale. The loss of granularity is partially offset by the explicit scaling factor $\alpha$ (the paper discusses values such as $128$ in this regime), giving a useful point on the optimization Pareto frontier: very small state, no SVD overhead, and still strong pre-training behavior.

\paragraph{Q-APOLLO.}
The APOLLO paper also stresses that the family combines naturally with \textbf{quantization} for ultra-low-memory training. This repository turns that systems observation into a concrete optimizer variant, \texttt{q\_apollo}. In the implementation, APOLLO's low-rank first and second moments are stored in quantized form together with per-tensor scales and offsets, and are dequantized only when needed for the next step. Q-APOLLO therefore preserves APOLLO's projected-gradient logic while reducing the precision of the remaining optimizer state, making it the most aggressive memory-saving member of the local optimizer stack.

\paragraph{Anon (Adaptivity Non-restricted Optimizer with Novel convergence technique).}
Anon \citep{anon2026anon} introduces continuously tunable adaptivity via a single parameter $\gamma \in \mathbb{R}$, bridging the gap between SGD-like ($\gamma=0$) and Adam-like ($\gamma=1$) behaviors, and extrapolating beyond both. The core innovation is the Incremental Delay Update (IDU) mechanism, which enables stable convergence across the entire real-valued adaptivity spectrum:

\begin{align}
m_t &= \beta_1 m_{t-1} + (1-\beta_1) g_t \quad\text{(first moment)}\\
v_t &= \beta_2 v_{t-1} + (1-\beta_2) g_t^2 \quad\text{(second moment)}\\
\bar{v}_t &= \frac{v_t}{1-\beta_2^{\max(t,1)}} + \epsilon \quad\text{(bias-corrected second moment)}\\
\text{fixed\_lr}_k &= \begin{cases}
\bar{v}_k^{-\gamma/2} & \text{if } k = 0 \\
\sqrt{2 / \left(\frac{1}{\text{fixed\_lr}_{k-1}^2} + \bar{v}_k^\gamma\right)} & \text{if } k > 0
\end{cases}\\
\theta_{t+1} &= \theta_t - \frac{\eta}{1-\beta_1^t} \cdot m_t \odot \text{fixed\_lr}_{\lfloor\log_2 t\rfloor}
\end{align}

Unlike standard adaptive optimizers, IDU updates the preconditioner only at powers of two ($t = 2^k$), accumulating gradient statistics between updates and aggregating them recursively. This soft, multi-scale accumulator avoids the hard max-tracking of AMSGrad while provably stabilizing convergence for any $\gamma \in \mathbb{R}$. The parameter $\gamma$ directly controls adaptivity: $\gamma > 1$ accelerates escape from saddle points (favorable for complex architectures), $\gamma = 1$ recovers Adam-like behavior, $\gamma = 0$ yields SGD-like scaling, and $\gamma < 0$ biases toward flatter minima (beneficial for classical architectures like CNNs).

Anon maintains three state buffers ($m$, $v$, fixed\_lr) with $\mathcal{O}(n)$ per-step cost, placing it in the same complexity class as AdamW but with the added flexibility of tunable adaptivity. The IDU technique provably converges at $\mathcal{O}(\sqrt{T})$ regret for convex problems and $\mathcal{O}(\ln T / \sqrt{T})$ for non-convex settings, matching the guarantees of mainstream adaptive optimizers while supporting the full real-valued adaptivity range.

\begin{algorithm}[H]
\caption{Anon (Adaptivity Non-restricted Optimizer with IDU)}
\label{alg:anon}
\begin{algorithmic}[1]
\Require Initial parameters $\theta_0$, learning rate $\eta$, $\beta_1, \beta_2$, tolerance $\epsilon$, adaptivity $\gamma \in \mathbb{R}$
\Ensure Updated parameters $\theta_T$
\State Initialize: $m \gets 0$, $v \gets 0$, fixed\_lr $\gets 0$, step $t \gets 0$, IDU counter $k \gets -1$
\For{$t = 1, 2, \ldots, T$}
    \State $g_t \gets \nabla f(\theta_{t-1})$
    \State $m \gets \beta_1 m + (1-\beta_1) g_t$
    \State $v \gets \beta_2 v + (1-\beta_2) g_t^2$
    \If{$t = 2^k$}
        \State $k \gets k+1$
        \State $\bar{v} \gets v/(1 - \beta_2^{\max(t/2,1)}) + \epsilon$
        \If{$k = 0$}
            \State fixed\_lr $\gets \bar{v}^{-\gamma/2}$
        \Else
            \State fixed\_lr $\gets \sqrt{2 / (1/\text{fixed\_lr}^2 + \bar{v}^\gamma)}$
        \EndIf
        \State $v \gets 0$
    \EndIf
    \State $\theta_t \gets \theta_{t-1} - \frac{\eta}{1-\beta_1^t} \cdot m \odot \text{fixed\_lr}$
\EndFor
\Return $\theta_T$
\end{algorithmic}
\end{algorithm}

\paragraph{Prodigy (Approximating the Distance Estimate).}
Prodigy adapts the effective step scale through a running distance-like statistic, eliminating the need for explicit learning-rate tuning. The algorithm maintains a cumulative distance estimate:
\begin{align}
u_t &= \frac{g_t}{\sqrt{v_t}+\epsilon} \quad\text{(unnormalized adaptive step)} \\
s_t &= s_{t-1} + \langle u_t, \theta_t - \theta_0 \rangle \quad\text{(cumulative signed distance)} \\
d_t &= \max(d_{t-1}, d_0 + d_{\text{coef}} \cdot s_t) \quad\text{(distance estimate with lower bound)} \\
\theta_{t+1} &= \theta_t - \eta \cdot d_t \cdot u_t
\end{align}
where $d_0$ is an initial bound and $d_{\text{coef}}$ controls how aggressively the estimate adapts. The distance estimate $d_t$ captures the combined magnitude of past gradients weighted by actual parameter displacement, implementing an implicit effective learning rate. This distance-aware scaling achieves robust convergence across problem scales and batch sizes without requiring a learning-rate schedule. Prodigy reduces hyperparameter sensitivity by estimating step sizes from optimization geometry rather than problem-specific prior knowledge.

\begin{algorithm}[H]
\caption{Prodigy (Approximating the Distance Estimate)}
\label{alg:prodigy}
\begin{algorithmic}[1]
\Require Initial parameters $\theta_0$, learning rate $\eta$, coeff $d_{\text{coef}}$, initial dist $d_0 > 0$
\Ensure Updated parameters $\theta_T$
\State Initialize: $s \gets 0$ (cumulative signed distance), $d \gets d_0$
\For{$t = 1, 2, \ldots, T$}
    \State $g_t \gets \nabla f(\theta_{t-1})$
    \State $v_t \gets \beta_2 v_{t-1} + (1-\beta_2) g_t^2$ \Comment{Second moment}
    \State $u_t \gets g_t / (\sqrt{v_t} + \epsilon)$ \Comment{Unnormalized adaptive}
    \State Update distance: $s \gets s + \langle u_t, \theta_t - \theta_0 \rangle$ \Comment{Signed distance}
    \State $d \gets \max(d, d_0 + d_{\text{coef}} \cdot s)$ \Comment{Lower-bounded distance}
    \State $\theta_t \gets \theta_{t-1} - \eta \cdot d \cdot u_t$ \Comment{Distance-scaled update}
\EndFor
\Return $\theta_T$
\end{algorithmic}
\end{algorithm}

\subsection{Second-Order, Geometric, and Orthogonality Optimizers}

\begin{algorithm}[H]
\caption{Shampoo (Matrix Preconditioning via Kronecker Products)}
\label{alg:shampoo}
\begin{algorithmic}[1]
\Require Initial parameters $\theta_0$, learning rate $\eta$, eigendecomposition interval $K$
\Ensure Updated parameters $\theta_T$
\State Initialize: $L \gets \epsilon I_m$, $R \gets \epsilon I_n$ (left and right Gram matrices)
\For{$t = 1, 2, \ldots, T$}
    \State $G \gets \nabla f(\theta_{t-1})$ \Comment{Gradient}
    \State Update Gram matrices: $L \gets L + G G^\top$, $R \gets R + G^\top G$
    \If{$t \mod K == 0$}
        \State $Q_L, \Lambda_L \gets \text{eigh}(L)$ \Comment{Eigendecomposition of $L$}
        \State $Q_R, \Lambda_R \gets \text{eigh}(R)$ \Comment{Eigendecomposition of $R$}
        \State $L^{-1/4} \gets Q_L (\Lambda_L + \epsilon)^{-1/4} Q_L^\top$
        \State $R^{-1/4} \gets Q_R (\Lambda_R + \epsilon)^{-1/4} Q_R^\top$
    \EndIf
    \State $\Delta \theta \gets L^{-1/4} G R^{-1/4}$ \Comment{Preconditioned gradient}
    \State $\theta_t \gets \theta_{t-1} - \eta \cdot \Delta \theta$
\EndFor
\Return $\theta_T$
\end{algorithmic}
\end{algorithm}

\paragraph{Shampoo (Matrix Preconditioning via Kronecker Products).}
Shampoo is a \textbf{structured second-order method} that computes matrix preconditioners from Kronecker-structured outer-product statistics. For a 2D parameter matrix $W \in \mathbb{R}^{m \times n}$:
\begin{align}
L_t &= L_{t-1} + G_t G_t^\top \quad\text{(left/row Gram matrix, $m \times m$)} \\
R_t &= R_{t-1} + G_t^\top G_t \quad\text{(right/column Gram matrix, $n \times n$)} \\
L_t^{-1/4} &= Q_L (\Lambda_L)^{-1/4} Q_L^\top \quad\text{(via eigendecomposition)} \\
R_t^{-1/4} &= Q_R (\Lambda_R)^{-1/4} Q_R^\top \\
\Delta W_t &= L_t^{-1/4} G_t R_t^{-1/4} \quad\text{(preconditioned gradient)}
\end{align}
Shampoo approximates the full-matrix Adagrad preconditioner $H^{-1/2}$ (where $H$ is the Hessian) using Kronecker-factored structure. Instead of storing a $(mn) \times (mn)$ preconditioner, Shampoo maintains two smaller matrices ($m \times m$ and $n \times n$), capturing cross-parameter correlations within and across groups. The method has proven effective at scale (Google production systems) and is well-suited to transformer architectures with high-rank structure. Eigendecompositions are performed periodically (e.g., every $K=10$ steps) to amortize cost. Tradeoff: higher per-step compute and periodic $O(m^3 + n^3)$ eigendecomposition versus  improved conditioning and accelerated convergence.

\paragraph{SOAP (Shampoo with Adam in eigenbasis).}
SOAP is a simplified variant of Shampoo that decouples preconditioning from momentum tracking. Instead of complex matrix algebra on preconditioned gradients, SOAP runs standard Adam in the eigenbasis of Shampoo's preconditioners:
\begin{align}
L_t &= L_{t-1} + G_t G_t^\top, \quad R_t = R_{t-1} + G_t^\top G_t \quad\text{(Gram accumulation)} \\
Q_L, \Lambda_L &= \text{eigh}(L_t), \quad Q_R, \Lambda_R = \text{eigh}(R_t) \quad\text{(periodic eigendecomposition)} \\
G_t^{\text{rot}} &= Q_L^\top G_t Q_R \quad\text{(rotate gradient into eigenbasis)} \\
m_t^{\text{rot}} &= \beta_1 m_{t-1}^{\text{rot}} + (1-\beta_1) G_t^{\text{rot}} \\
v_t^{\text{rot}} &= \beta_2 v_{t-1}^{\text{rot}} + (1-\beta_2)(G_t^{\text{rot}})^2 \quad\text{(Adam in rotated space)} \\
U_t^{\text{rot}} &= \frac{m_t^{\text{rot}}}{\sqrt{v_t^{\text{rot}}}+\epsilon}, \quad U_t = Q_L U_t^{\text{rot}} Q_R^\top \quad\text{(rotate back)}
\end{align}
The key insight: all momentum tracking occurs in the well-conditioned eigenbasis, simplifying numerical stability and theoretical analysis. SOAP combines benefits of Shampoo (explicit curvature structure) and Adam (proven momentum mechanics), while being cleaner and often more stable than full Shampoo. Eigendecompositions are recomputed every $K$ steps, amortizing the cost.

\begin{algorithm}[H]
\caption{SOAP (Shampoo with Adam in Eigenbasis)}
\label{alg:soap}
\begin{algorithmic}[1]
\Require Initial parameters $\theta_0$, learning rate $\eta$, $\beta_1, \beta_2$, eigendecomposition interval $K$
\Ensure Updated parameters $\theta_T$
\State Initialize: $L \gets \epsilon I_m$, $R \gets \epsilon I_n$, $m \gets 0$, $v \gets 0$
\For{$t = 1, 2, \ldots, T$}
    \State $G \gets \nabla f(\theta_{t-1})$ \Comment{Gradient}
    \State Update Gram: $L \gets L + G G^\top$, $R \gets R + G^\top G$
    \If{$t \mod K == 0$}
        \State $Q_L, \Lambda_L \gets \text{eigh}(L)$, $Q_R, \Lambda_R \gets \text{eigh}(R)$
    \EndIf
    \State $G^{\text{rot}} \gets Q_L^\top G Q_R$ \Comment{Rotate gradient into eigenbasis}
    \State $m \gets \beta_1 m + (1-\beta_1) G^{\text{rot}}$ \Comment{Exponential moving average (bias-correct outside)}
    \State $v \gets \beta_2 v + (1-\beta_2) (G^{\text{rot}})^2$ \Comment{Second moment (bias-correct outside)}
    \State $u \gets m / (\sqrt{v} + \epsilon)$ \Comment{Adaptive step in eigenbasis}
    \State $\Delta \theta \gets Q_L u Q_R^\top$ \Comment{Rotate back to parameter space}
    \State $\theta_t \gets \theta_{t-1} - \eta \cdot \Delta \theta$
\EndFor
\Return $\theta_T$
\end{algorithmic}
\end{algorithm}

\paragraph{Lion (EvoLved Sign Momentum).}
Lion achieves \textbf{minimal memory overhead} by using sign-based updates instead of full adaptive scaling:\begin{align}
c_t &= \beta_1 m_t + (1-\beta_1) g_t \quad\text{(momentum input)} \\
\theta_{t+1} &= \theta_t - \eta \left(\text{sign}(c_t) + \lambda\theta_t\right) \quad\text{(sign operator update)} \\
m_{t+1} &= \beta_2 m_t + (1-\beta_2) g_t \quad\text{(momentum accumulation)}
\end{align}
All element-wise scaling operations are replaced with \texttt{sign()}, which outputs $\{-1, 0, +1\}$. This dramatically reduces memory compared to Adam-like methods: Lion stores only the momentum buffer, no second moment variance. The tradeoff: sign-based updates sacrifice the element-wise learning-rate adaptation that makes Adam effective. Lion is positioned not as a universally superior optimizer but as a specialized low-memory, high-throughput alternative for scenarios where activation memory dominates and some optimizer sophistication can be sacrificed. Works well in large-batch regimes and when hardware throughput is the primary constraint.

\begin{algorithm}[H]
\caption{Lion (EvoLved Sign Momentum)}
\label{alg:lion}
\begin{algorithmic}[1]
\Require Initial parameters $\theta_0$, learning rate $\eta$, $\beta_1, \beta_2$, weight decay $\lambda$
\Ensure Updated parameters $\theta_T$
\State Initialize: $m \gets 0$ (momentum buffer)
\For{$t = 1, 2, \ldots, T$}
    \State $g \gets \nabla f(\theta_{t-1})$ \Comment{Gradient}
    \State $c \gets \beta_1 m + (1-\beta_1) g$ \Comment{Momentum input}
    \State $\theta_t \gets \theta_{t-1} - \eta (\text{sign}(c) + \lambda \theta_{t-1})$ \Comment{Sign-based update with weight decay}
    \State $m \gets \beta_2 m + (1-\beta_2) g$ \Comment{Momentum accumulation (decoupled)}
\EndFor
\Return $\theta_T$
\end{algorithmic}
\end{algorithm}

\paragraph{Sophia (Second-Order Hessian Information with Optimized Approximation).}
Sophia uses \textbf{diagonal Hessian estimates} for curvature-aware scaling without the memory and compute cost of dense second-order methods.\begin{align}
m_t &= \beta_1 m_{t-1} + (1-\beta_1) g_t \quad\text{(first moment)} \\
h_t &= \beta_2 h_{t-(k-1)} + (1-\beta_2) \hat{h}_t \quad\text{(diagonal Hessian, updated every $k$ steps)} \\
\text{clip}(x, C) &= \min(\max(x, -C), C) \quad\text{(element-wise clipping)} \\
\theta_{t+1} &= \theta_t - \eta \cdot \text{clip}\left(\frac{m_t}{\max(\gamma h_t, \epsilon)}, 1\right)
\end{align}
where $\hat{h}_t$ is a diagonal estimate (e.g., Hutchinson-trace estimator) and $\gamma$ is a scaling factor. The diagonal Hessian $h_t$ captures local curvature, allowing the optimizer to take smaller steps in sharp directions and larger steps in flat directions. The clipping operation $\text{clip}(\\cdot, 1)$ prevents adaptive steps from exploding. Sophia belongs to the family of curvature-aware methods seeking better conditioning without the $O(n^2)$ or $O(n^3)$ cost of full second-order methods. Works particularly well in second-pass fine-tuning scenarios.

\begin{algorithm}[H]
\caption{Sophia (Diagonal Hessian with Clipped Updates)}
\label{alg:sophia}
\begin{algorithmic}[1]
\Require Initial parameters $\theta_0$, learning rate $\eta$, $\beta_1, \beta_2$, Hessian update freq $k$, clip threshold $C$ 
\Ensure Updated parameters $\theta_T$
\State Initialize: $m \gets 0$, $h \gets \epsilon$ (diagonal Hessian estimate)
\For{$t = 1, 2, \ldots, T$}
    \State $g \gets \nabla f(\theta_{t-1})$
    \State $m \gets \beta_1 m + (1-\beta_1) g$ \Comment{First moment (bias-correct outside)}
    \If{$t \mod k == 0$}
        \State $\hat{h} \gets \text{HutchinsonEstimate}(g)$ \Comment{Diagonal Hessian via Hutchinson}
        \State $h \gets \beta_2 h + (1-\beta_2) \hat{h}$ \Comment{Update Hessian estimate}
    \EndIf
    \State $u \gets \frac{m}{\max(\gamma h, \epsilon)}$ \Comment{Adaptive step (element-wise)}
    \State $u \gets \text{clip}(u, C)$ \Comment{Element-wise clipping to $[-C, C]$}
    \State $\theta_t \gets \theta_{t-1} - \eta \cdot u$
\EndFor
\Return $\theta_T$
\end{algorithmic}
\end{algorithm}

\paragraph{Muon and Turbo-Muon (Orthogonality-Based Optimizers).}
Muon and Turbo-Muon are \textbf{orthogonality-oriented optimizers} that reshape update geometry using Newton-Schulz polynomial iterations for orthogonalization. For a 2D parameter matrix $W$ with gradient $G_t$:
\begin{align}
X_0 &= \frac{G_t}{\|G_t\|_F + \epsilon} \quad\text{(normalized gradient)} \\
A_k &= X_k X_k^\top \quad\text{(Gramian)} \\
B_k &= b A_k + c A_k^2 \quad\text{(polynomial step, coefficients $b, c$ from Newton-Schulz)} \\
X_{k+1} &= a X_k + B_k X_k \quad\text{(iteration $k=0,1,\ldots,4$)} \\
W_{t+1} &= W_t - \eta X_K \quad\text{(update with orthogonalized direction)}
\end{align}
Muon performs 5 Newton-Schulz iterations to produce an approximately orthogonal direction, ensuring updates respect geometric constraints. Turbo-Muon adds an \textbf{almost-orthogonal preconditioning} (AOL) step before iterations, reducing the number of required orthogonalization steps to 4. The rationale: orthogonal updates preserve norms during training, avoiding the norm creep observed in momentum-based methods. These optimizers show promise on large-scale training but require custom CUDA kernels for efficiency. Tradeoff: high per-step compute (matrix multiplications and polynomial iterations) versus fundamentally better-conditioned update geometry.

\begin{algorithm}[H]
\caption{Muon (Newton-Schulz Orthogonalization)}
\label{alg:muon}
\begin{algorithmic}[1]
\Require Initial parameters $\theta$ (matrices), learning rate $\eta$, Newton-Schulz iters $K$
\Ensure Updated parameters $\theta$
\For{each 2D parameter matrix $W$}
    \State $G \gets \nabla f(W)$ \Comment{Gradient}
    \State $X_0 \gets \frac{G}{\|G\|_F + \epsilon}$ \Comment{Normalize gradient by Frobenius norm}
    \For{$k = 0, 1, \ldots, K-1$}
        \State $A_k \gets X_k X_k^\top$ \Comment{Gramian (cost: $O(mn^2)$ or $O(m^2n)$ for tall/wide)}
        \State $B_k \gets (3/2) A_k - (1/2) A_k^2$ \Comment{Newton-Schulz coefficients}
        \State $X_{k+1} \gets B_k X_k$ \Comment{Polynomial step}
    \EndFor
    \State $W \gets W - \eta X_K$ \Comment{Update with orthogonalized direction}
\EndFor
\Return updated $\theta$
\end{algorithmic}
\end{algorithm}

\begin{algorithm}[H]
\caption{Turbo-Muon (AOL-Preconditioned Orthogonalization)}
\label{alg:turbo-muon}
\begin{algorithmic}[1]
\Require Initial parameters $\theta$ (matrices), learning rate $\eta$, Newton-Schulz iters $K'$ (typically 4)
\Ensure Updated parameters $\theta$
\For{each 2D parameter matrix $W$}
    \State $G \gets \nabla f(W)$ \Comment{Gradient}
    \State $X_0 \gets \frac{G}{\|G\|_F + \epsilon}$ \Comment{Normalize gradient}
    \State \bf{AOL (Almost-Orthogonal preconditioner):} \Comment{Preconditioning step}
    \State $P \gets (1.5 I - 0.5 X_0 X_0^\top)$ \Comment{Preconditioning matrix}
    \State $X_0 \gets P X_0$ \Comment{Preconditioned start}
    \For{$k = 0, 1, \ldots, K'-1$}
        \State $A_k \gets X_k X_k^\top$
        \State $B_k \gets (3/2) A_k - (1/2) A_k^2$
        \State $X_{k+1} \gets B_k X_k$ \Comment{Reduced iterations due to AOL preconditioning}
    \EndFor
    \State $W \gets W - \eta X_{K'}$ \Comment{Update with preconditioned orthogonalized direction}
\EndFor
\Return updated $\theta$
\end{algorithmic}
\end{algorithm}

\small
\begin{longtable}{p{2.5cm}p{3.5cm}p{2.0cm}p{5.3cm}}
\toprule
\textbf{Group} & \textbf{Methods} & \textbf{Primary Goal} & \textbf{Interpretation from the Survey} \\
\midrule
\endhead
Classical baseline & SGD, AdamW, RAdam & stability and reference baselines & These define the comparison floor for newer optimizer claims. \\
Momentum redesign & Adan, AdEMAMix, MARS, Cautious AdamW & faster or safer first-order adaptation & Best when convergence speed or noisy-gradient stability is the main concern. \\
Large-batch and schedule simplification & LAMB, Schedule-Free AdamW & operational robustness at scale & Reduce brittleness from batch-size growth or schedule engineering. \\
Memory-efficient & Adafactor, GaLore, APOLLO, APOLLO-Mini, Q-APOLLO, Lion & optimizer-state reduction & Most useful when VRAM is dominated by optimizer state rather than activations; APOLLO-family methods replace dense AdamW moments with projected or quantized structured scaling. \\
Curvature-aware & Shampoo, SOAP, Sophia & better conditioning & Prefer when richer geometry is worth implementation and compute overhead. \\
Geometry-oriented & Muon, Turbo-Muon & orthogonalized update structure & Specialized options for matrix geometry and representation shaping. \\
\bottomrule
\end{longtable}
\normalsize

